% !PRINTED 104
\weberpage{104}
Zur Bestimmung der Konstanten $A$ wenden wir, wie in der Folge häufig, das Mittel an, daß wir $u = 0$ setzen, in (1) nach der Differentiation, und dann rechts und links von der Formel (5), § 23
\weq{5}{\vartheta'_{11} = \pi\,\vartheta_{00}\,\vartheta_{10}\,\vartheta_{01}}
Gebrauch machen; so folgt
\[
A^2 = e^{\frac{\pi i}{2}}, \quad A = e^{\frac{\pi i}{4}};
\]
daß bei $A$ das positive Zeichen steht, ergibt sich aus einer der Formeln (3), (4) nach der Schlußbemerkung von § 25, wenn man $\omega$ unendlich werden läßt.

Sonach erhält man
\weq{6}{\begin{aligned}
\vartheta_{11}(u, \omega+1) &= e^{\frac{\pi i}{4}}\vartheta_{11}(u) \\
\vartheta_{10}(u, \omega+1) &= e^{\frac{\pi i}{4}}\vartheta_{10}(u) \\
\vartheta_{01}(u, \omega+1) &= \vartheta_{00}(u) \\
\vartheta_{00}(u, \omega+1) &= \vartheta_{01}(u).
\end{aligned}}
II. $\begin{pmatrix} 0, & 1 \\ -1, & 0 \end{pmatrix}$ oder $\left(\omega, \dfrac{-1}{\omega}\right)$.

Es ist wieder nach § 27, (12)
\weq{7}{e^{-\frac{\pi i u^2}{\omega}}\vartheta_{11}\left(\frac{u}{\omega}, -\frac{1}{\omega}\right) = A\vartheta_{11}(u, \omega),}
und durch Vermehrung von $u$ um $\dfrac{1}{2}$, $\dfrac{\omega}{2}$, $\dfrac{1+\omega}{2}$
\weq{8}{e^{-\frac{\pi i u^2}{\omega}}\vartheta_{01}\left(\frac{u}{\omega}, -\frac{1}{\omega}\right) = iA\vartheta_{10}(u, \omega),}
\weq{9}{e^{-\frac{\pi i u^2}{\omega}}\vartheta_{10}\left(\frac{u}{\omega}, -\frac{1}{\omega}\right) = iA\vartheta_{01}(u, \omega),}
\weq{10}{e^{-\frac{\pi i u^2}{\omega}}\vartheta_{00}\left(\frac{u}{\omega}, -\frac{1}{\omega}\right) = iA\vartheta_{00}(u, \omega),}
woraus man wie oben erhält:
\[
A = \pm i\sqrt{-i\omega}.
\]
Aus $u = 0$ und einem rein imaginären $\omega$ schließt man, daß, wenn $\sqrt{-i\omega}$ so genommen wird, daß \gesperrt{der reelle Teil positiv ist, das untere Zeichen stehen muß}, und es ergibt sich daher:
%CONT>>
